\section{导读：为什么这期是 Agent 的技术谱系课}

这期不是普通嘉宾访谈，而是一堂口头版 Agent 技术史。嘉宾苏煜是俄亥俄州立大学计算机系教授、NeoCognition 创始人，长期研究 Language Agent、Computer Use Agent、Mind2Web、LM Planner、MMMU 等方向。访谈试图回答一个大问题：为什么 2026 年的 AI 叙事从 Chat 进入 Agent，为什么 OpenClaw 这样的产品让行业产生类似 ChatGPT Moment 的震动，以及为什么 browser、desktop、mobile、GUI、CLI、API、coding 这些边界正在被重新组织。

\begin{importantbox}{本期的核心问题}
如果把 ChatGPT Moment 看成 LLM 范式的公众确认，那么 OpenClaw Moment 更像是 Agent 范式的公众确认：模型不再只是回答问题，而是进入环境、使用工具、执行动作、持续学习，并逐渐逼近 universal digital agent。
\end{importantbox}




\subsection{本章小结}

本期适合被整理成“Agent 技术谱系”笔记：先定义 Agent，再回顾 logical agent、neural agent、semantic parsing、language agent，最后讨论 OpenClaw、universal digital agent、continual learning、world model 与产业扩散。

